% Sección: aprendizaje por diferencia temporal (TD).
\section{Aprendizaje por Diferencia Temporal}
\label{sec:td}

El primer método de entrenamiento aprende la función de valor por
\emph{bootstrapping}, sin esperar al final del episodio, sobre la misma
arquitectura lineal $\hat{U}(\sigma; w) = w^{\top}\varphi(\sigma)$ de la
Sección~\ref{sec:valor}. Es el baseline funcional del proyecto.

\subsection{Error TD y actualización semi-gradiente}
Para una transición entre afterstates $\sigma_t \to \sigma_{t+1}$ con recompensa
$R_{t+1}$, el objetivo TD(0) y el error TD son
\begin{align}
  y_t &= R_{t+1} + \gamma\, \hat{U}(\sigma_{t+1}; w), \label{eq:td-target}\\
  \delta_t &= y_t - \hat{U}(\sigma_t; w)
           = R_{t+1} + \gamma\, w^{\top}\varphi(\sigma_{t+1})
             - w^{\top}\varphi(\sigma_t). \label{eq:td-error}
\end{align}
La actualización \emph{semi-gradiente} desciende sobre $\tfrac{1}{2}\delta_t^{2}$
tratando el objetivo $y_t$ como constante (no se deriva respecto a $w$):
\begin{equation}
  w \leftarrow w + \alpha\, \delta_t\, \nabla_w \hat{U}(\sigma_t; w)
            = w + \alpha\, \delta_t\, \varphi(\sigma_t),
  \label{eq:td-update}
\end{equation}
con tasa de aprendizaje $\alpha$. En el estado terminal (top-out) se toma
$\hat{U} = 0$, de modo que $\delta_T = R_T - w^{\top}\varphi(\sigma_T)$.

\subsection{Política de comportamiento}
Durante el entrenamiento se sigue una política $\varepsilon$-greedy sobre la
acción greedy de \eqref{eq:politica}: con probabilidad $1-\varepsilon$ se elige
$\arg\max_a w^{\top}\varphi(f(s,a))$ y con probabilidad $\varepsilon$ una
colocación al azar (exploración). Las tasas $\alpha$ y $\varepsilon$ decaen con
los episodios. El descuento $\gamma < 1$ (en la práctica $\gamma \in \{0.99, 0.999\}$,
siendo $0.999$ el de la mejor configuración) actúa como recurso algorítmico que
estabiliza el \emph{bootstrapping}; el desempeño que se reporta sigue siendo el
total de líneas por partida.

\subsection{Algoritmo}
\begin{verbatim}
Inicializar w = 0
para cada episodio:
    s = reset();  a = epsilon_greedy(legal_placements(s))
    mientras no termine:
        sigma = afterstate(s, a)
        (s', R, fin) = step(a)
        a' = epsilon_greedy(legal_placements(s'))      # None si no hay jugada
        sigma' = afterstate(s', a')                    # terminal si a' = None
        delta = R + gamma * U(sigma') - U(sigma)       # U(terminal) = 0
        w = w + alpha * delta * phi(sigma)
        s, a = s', a'
\end{verbatim}
El entrenamiento se paraleliza de forma asíncrona: varios trabajadores juegan
episodios y acumulan sus actualizaciones sobre un vector de pesos en memoria
compartida (esquema tipo Hogwild!).

\subsection{Estabilidad}
La actualización \eqref{eq:td-update} es \emph{semi-gradiente}: no es el gradiente
de una función objetivo fija, por lo que la combinación de \emph{bootstrapping},
aproximación de funciones y muestreo puede hacer \textbf{divergir} los pesos
\citep{sutton2018}. En la práctica observamos que $\lVert w \rVert$ crece sin
acotarse, aunque la política greedy sigue rindiendo gracias a la invarianza de
escala (Sección~\ref{sec:valor}). Esta limitación motiva el método sin
\emph{bootstrap} ni problema de escala de la Sección~\ref{sec:cem}.
